\section{2 июня 1926 г.}

\lp{20260606_184836104}

Москва, 2 июня 1926 г.

Дорогой и милый мой Коля, ты, конечно, уже знаешь от Льва Эд\rem*{уардовича} о
скоропостижной кончине Егора Львовича Катуара. Смерть его своей неожиданностью
глубоко поразила всех знавших и любивших его. Для меня, прожившего в дружбе с
ним почти всю жизнь (я познакомился с ним еще будучи гимназистом), смерть его
явилась тяжелым испытанием, потерей прекрасного и благороднейшего человека,
когда круг близких все редеет и редеет -- очень тяжело. Последний раз я навещал
его за три дня до его неожиданной смерти; он был как всегда приветлив, ласков и
совершенно бодр; говорил, что чувствует себя совсем хорошо (у него перед этим
было легкое недомогание). Накануне смерти, он весь день провел в консерватории
и был совершенно здоров; вечером лег спать как всегда в 11 часов; уснул -- и
больше не проснулся. Врач определил причину смерти как приступ грудной жабы.
Странно, что Егор Львович никогда на сердце не жаловался.

\lp{20260606_184850504}

Теперь нашим долгом является помочь его детям и внукам, которых он так любил и
за судьбу которых всегда очень беспокоился, и принять все меры к ознакомлению
не только русского общества, но и заграничного с его произведениями, столь
несправедливо оставшимися в тени и известными лишь весьма ограниченному числу
музыкантов. Я хорошо знаю как глубоко страдал от того Егор Львович; он часто
говорил мне, что самое тяжелое в жизни -- это сознание, что вся твоя работа,
все твои произведения никому не нужны. Я всячески старался его уверить в
неправильности его взгляда, но каждый раз безуспешно. И, конечно, он имел
основания говорить, что его, как композитора, не хотят даже замечать.

В Москве мы постараемся, сделаем все возможное, чтобы в ближайший сезон
устроить целый ряд концертов из его произведений. (Александр Борисович --
первый, который устраивает теперь же, 8 июня. Будут исполнять трио,
фортепианный квинтет и романсы), но было бы черезвычайно важно устроить

\lp{20260606_184850504}

\noindent ряд концертов за границей, и я очень надеюсь, что ты, С. В.
\rem{Рахманинов} и Лев Эдуардович \rem{Конюс} в этом нам поможете. Как это
лучше сделать я не знаю, но думаю, что для начала, если Лев, ты и Рахманинов
написали бы какую-то небольшую статью-некролог в заграничных газетах, главным
образом в американских, немецких и английских -- то это было бы очень хорошо.

Теперь относительно детей Егора Львовича; Петя находится в Берлине, часто
выступает и живет самостоятельно. Ему надо бы оказать поддержку, главным
образом в его выступлениях. Он мог бы отлично сыграть сонаты и камерные
сочинения Егора Львовича; адреса его я не знаю, но Лев Эдуардович наверное
знает, т.к. он теперьпороднился с Катуарами. Теперь относительно дочерей: Аня,
Лена и Таня; живут, или вернее сказать, жили с отцом (они остались в той же
квартире, уступив после смерти отца одну комнату).

\lp{20260606_184857012}

Старшая, Аня, замужем за Яковенко (певец) жила со своими двумя маленькими
детьми всецело на иждивении отца, т.к. ее муж почти ничего не зарабатывает (он
прекрасный и чистый человек, но совсем не от мира сего и семью он содержать не
будет в силах). Лена, замужем за Немировичем (доцент высшего учебного
заведения), эти жили самодостаточно в материальном отношении, и, наконец,
младшая Таня, незамужем, учится в консерватории (по фортепиано у меня в
классе). Окончит курс не раньше чем через два года. Жила на иждивении отца.

Александр Борисович, подал прошение Луначарскому о назначении Ане пенсии, но
когда это выйдет неизвестно. Тане мы постараемся выхлопотать стипендию рублей в
20 в месяц в консерватории. Это можно будут сделать не раньше начала учебного
года. Заговорив о детях Егора Львовича, мне пришло в голову следущее: в конце
апреля я написал письма тебе и Сергею Васильевичу по твоему адресу.

\lp{20260606_184956770}

В письме к Сергею Васильевичу я между прочим писал, что мне крайне тяжело
получать от него денежную помощь (он мне посылал иногда через месяц, а иногда
каждый месяц по 10 долларов), и я в этом письме благодарю его за внимание,
очень прошу денег мне не посылать, т.к. пока я здоров и работаю, мне просто это
тяжело. Не получив ответа ни от тебя, ни от Сергея Васильевича, я на днях опять
получил в конторе банка переводом 10 долларов. Милый Коля, попроси
незамедлительно С.В., если он мое письмо не получил, не посылать мне больше
денег; передай ему мою сердечную благодарность за продовольственные посылки в
голодные годы и скажи, что если бы он имел возможность и желание, то пусть он
ежемесячно посылает эти 10 долларов Тане или Ане. Адрес их: Петровский бульвар,
д.~8 Татьяна Георгиевна Катуар или Анна Георгиевна Яковенко. Пусть он это
сделает вместо венка на могилу милого и дорогого Егора Львовича.

\lp{20260606_185010396}

Милый мой и хороший Коля, теперь еще один важнейший вопрос в связи со смертью
Егора Львовича. -- Жизнь не останавливается после смерти людей, и, не успели мы
схоронить Егора Львовича, как поднялся вопрос о кандидатах на его место. Мне не
хотелось бы называть этих кандидатов, скажу только, что мне кажется, что если
ты скучаешь по России и мечтаешь о возвращении в Москву, то ты явился бы для
нас единственным кандидатом на это место во всех отношениях. К сожалению «для
нас» не означает всей консерватории. Вот в этом то весь вопрос, и поэтому то я
и пишу тебе. Бацилла, которая называется в современной музыке «искание новых
путей», завела в тупик не только учащихся (т.к. учащиеся лишь спешат подражать
старшим), но и часть профессуры по этому отделу (творческий факультет, как он
теперь называется), и эти «искания» много крови стоили покойному Егору
Львовичу. Связь с музыкой порывается, о гармонии

\lp{20260606_185010396}

\noindent
не имеют больше никакого представления, хотя бы в самом примитивном смысле. Все
это направление пришло к нам с запада, пустило глубочайшие корни, затопило и
захватило почти всю молодежь. Надо к этому для ясности прибавить, что это ни в
какой связи с политикой не стоит, как раз наоборот. В Москве имеется некоторый
ряд композиторов-коммунистов, которые ведут с этим направлением борьбу, но они
слишком слабы, как музыканты, чтобы иметь влияние на молодежь. К сожалению, и
более сильные музыканты имеют все меньше и меньше влияния в этом отношении.
Думается мне, что если бы ты собрался в Москву вместо февраля, как ты
предполагал, в сентябре, и попробовал бы взяться за дело со свежими силами, то
кто знает, может быть ты, пробыв хотя бы 1 - 2 года для опыта, сумел бы оказать
влияние хотя бы на небольшую часть консерваторской молодежи и вернул бы их на
правильный путь и дал бы им возможность вновь научиться понимать и любить
музыку.

\lp{20260606_185021754}

\noindent
Дело это, конечно, весьма серьезное, и я ни в какой мере не берусь тебе давать
советы, т.к. этот призыв в Москву вызовет полную пертурбацию во всей твоей
жизни, но написать тебе об этом я счел своим долгом.

Напиши мне, получил ли ты ноты, которые я тебе послал через музыкальный сектор?

Передай мой сердечный поклон, а также и от Екатерины Петровны, Анне Михайловне,
Льву Эдуардовичу, Сергею Васильевичу, Ольге Николаевне и всем тем, кто нас
помнит. Вчера, первый раз за целый год, я был с Пантелеймоном Ивановичем в
Большом Театре на «Китеже». Ушел из театра глубоко взволнованный величавым и
прекрасным произведением. Образы Февроньи и Гришки захватили меня до глубины
души. Поставлена опера совершенно исключительно хорошо. С.К. на высоте; а сцены
в б. \rem{балете} «Китеже», как зрелище -- прямо потрясающе хороши.

Крепко-крепко тебя целую, мой милый Коля. Пришли мне, как только выйдет у тебя,
твою скрипичную сонату.

Твой А. Гедике.

\lp{20260606_184836104}

Сообщи мне свой летний адрес и адрес Сергея Васильевича и Льва Эдуардовича.
